% ============================================================
% Satterthwaite-Box Corrected J Test for GMM with Many
% Correlated Moment Conditions
%
% Required packages: amsmath, amssymb, booktabs, tcolorbox,
%                    geometry, hyperref, enumitem
% ============================================================

\section{A Corrected Overidentification Test via the
         Satterthwaite--Box Approximation}
\label{sec:jtest}

% ── Why we need this section ────────────────────────────────────────────────
The standard Hansen (1982) overidentification test --- the $J$ test ---
checks whether the moment conditions imposed by our GMM estimator are
consistent with the data.  Intuitively it asks: after choosing $\hat\beta$
to make the moments as close to zero as possible, are the
\emph{remaining} deviations small enough to be consistent with sampling
noise?  If yes, we do not reject the model; if the deviations are too
large, something is mis-specified.

In our setting two features of the data make the standard $J$ test break
down.  The goal of this section is to explain, step by step, what those
problems are and exactly how the Satterthwaite--Box correction fixes them.
We aim to make every step accessible to a reader who has seen basic
statistics but not necessarily advanced econometrics.

% ─────────────────────────────────────────────────────────────────────────────
\subsection{Background: What the Standard J Test Does}
\label{subsec:jtest_standard}
% ─────────────────────────────────────────────────────────────────────────────

Suppose we have $m$ moment conditions stacked into a vector
$\Delta \in \mathbb{R}^m$ and $k$ parameters $\beta \in \mathbb{R}^k$
with $m > k$ (more conditions than parameters, i.e.\ the model is
\emph{overidentified}).  The moment conditions say
\begin{equation}
  \mathbb{E}[\Delta - Q_H\beta] = \mathbf{0},
  \label{eq:moments}
\end{equation}
where $Q_H$ is the $m\times k$ matrix that maps parameters to their
predicted moment values.  At the true $\beta_0$ all $m$ conditions
should hold on average.

The GMM estimator picks $\hat\beta$ to make $\hat g \;{:=}\; \Delta - Q_H\hat\beta$
as close to zero as possible in a weighted sense.  The residual vector
$\hat g$ captures ``how well'' the model fits.  Ideally it should be
small --- but not too small.

The standard $J$ statistic is
\begin{equation}
  J \;=\; n\cdot \hat g'\,\hat\Omega^{-1}\,\hat g,
  \label{eq:jstandard}
\end{equation}
where $n$ is the sample size and $\hat\Omega$ is an estimate of the
covariance matrix of $\Delta$.  Under the null hypothesis that the model
is correctly specified, $J \;\overset{d}{\longrightarrow}\; \chi^2(m - k)$.
The degrees of freedom $m-k$ count the number of overidentifying
restrictions --- the number of moment conditions left over after
identifying all $k$ parameters.

\medskip
\noindent\textbf{When does this work?}  The derivation of $J\sim\chi^2(m-k)$
requires two things: (i)~$\hat\Omega$ is invertible, and (ii)~$\hat\beta$
is estimated using the same $\hat\Omega^{-1}$ as the weighting matrix
(efficient GMM).  Both conditions fail in our setting.

% ─────────────────────────────────────────────────────────────────────────────
\subsection{Problem 1: The Covariance Matrix is Singular}
\label{subsec:singularity}
% ─────────────────────────────────────────────────────────────────────────────

\paragraph{What singularity means.}
A matrix is \emph{singular} if it has no inverse.  Think of a $2\times2$
matrix whose two rows are identical: you cannot invert it because the
two rows carry the same information.  More generally, if a matrix has
fewer independent directions than it has rows, it is singular and
$\hat\Omega^{-1}$ in~\eqref{eq:jstandard} cannot be computed.

\paragraph{Why $\hat\Omega$ is singular here.}
Every one of our $m$ DiD moment conditions has the form
\begin{equation}
  \Delta_{g,c,s,t}
    = \bar Y_g(t) - \bar Y_g(s) - \bar Y_c(t) + \bar Y_c(s),
  \label{eq:did_form}
\end{equation}
where $\bar Y_g(t)$ is the average outcome of cohort $g$ at time $t$.
All $m$ moments are constructed from the same small set of
\emph{cohort-time averages} $\{\bar Y_g(t)\}$.  If there are $G$
cohorts and $T$ time periods, there are at most $G\times T$ distinct
averages --- yet $m$ can be thousands.

\medskip
\begin{tcolorbox}[colback=gray!8, colframe=gray!40, title=Simple analogy]
Suppose you have three numbers: $A$, $B$, $C$.
You compute six ``moments'':
\[
  A-B,\quad B-C,\quad A-C,\quad 2A-B-C,\quad A+B-2C,\quad A-2B+C.
\]
No matter how many moments you compute, they all live in a
three-dimensional space (the span of $A$, $B$, $C$).
Their $6\times6$ covariance matrix has rank at most 3.
Inverting a $6\times6$ matrix of rank 3 is impossible.
\end{tcolorbox}
\medskip

In our application with $m = 5{,}829$ DiD moments built from roughly
$G\times T \approx 1{,}000$ cohort-time averages, the covariance matrix
$\hat\Omega$ has rank at most ${\sim}1{,}000$.  Computing
$\hat\Omega^{-1}$ directly is impossible.

% ─────────────────────────────────────────────────────────────────────────────
\subsection{Problem 2: Downward Bias from Full-Sample Estimation}
\label{subsec:downward_bias}
% ─────────────────────────────────────────────────────────────────────────────

\paragraph{What the bias is.}
In efficient GMM, $\hat\beta$ is chosen by minimising
$(\Delta - Q_H\beta)'\hat\Omega^{-1}(\Delta - Q_H\beta)$.
The first-order conditions are
\begin{equation}
  Q_H'\,\hat\Omega^{-1}\,\hat g \;=\; \mathbf{0}.
  \label{eq:foc}
\end{equation}
This says the residuals $\hat g$ are \emph{forced} to be orthogonal
(in the $\hat\Omega^{-1}$ metric) to the columns of $Q_H$.
As a result $\hat g$ is as small as possible in the weighted norm ---
the estimator ``uses up'' the data to shrink the residuals.

When you then compute $J = n\cdot\hat g'\hat\Omega^{-1}\hat g$, you are
measuring the length of residuals that have already been minimised.
They are smaller than they would be at the true $\beta_0$, so
$\mathbb{E}[J] < m - k$ in finite samples.  The $J$ statistic is
\emph{biased downward}: it makes the model look better-fitting than it
really is, leading to under-rejection.

\medskip
\begin{tcolorbox}[colback=gray!8, colframe=gray!40,
  title=Simple analogy]
Imagine throwing darts at a board and measuring the average miss
distance.  If someone is \emph{allowed to move the board after seeing
where the darts land}, the board will always end up close to the darts
--- the measured ``miss'' is artificially small.  Using $\hat\beta$ from
the full set of moments is like moving the board: the residuals look
small because the estimator was specifically chosen to make them small.
\end{tcolorbox}

% ─────────────────────────────────────────────────────────────────────────────
\subsection{The Corrected Procedure, Step by Step}
\label{subsec:procedure}
% ─────────────────────────────────────────────────────────────────────────────

We now describe the full corrected procedure.  Each step is paired with
an explanation of which problem it addresses.

% ── Step 1 ──────────────────────────────────────────────────────────────────
\paragraph{Step 1: Estimate $\hat\beta$ from a \emph{restricted} set of
moments.}

Choose a subset of $k$ moment conditions --- just enough to identify
all $k$ parameters.  This is called the \emph{just-identified} or
\emph{restricted} set.  Solve the $k$ equations
\begin{equation}
  \Delta_{\text{res}} - Q_{H,\text{res}}\,\hat\beta_r \;=\; \mathbf{0}
  \label{eq:restricted_est}
\end{equation}
exactly for $\hat\beta_r$, where the subscript ``res'' denotes the
chosen subset.

\medskip
\noindent\textbf{Why this fixes Problem 2.}  The restricted estimator
$\hat\beta_r$ satisfies the first-order conditions
\eqref{eq:foc} \emph{only for the $k$ restricted moments}, not for the
remaining $m - k$ overidentifying moments.  Those overidentifying
residuals are free to reflect genuine specification error.  We are no
longer moving the board after seeing all the darts --- we fix the board
using a small held-out subset of darts and check the rest honestly.

% ── Step 2 ──────────────────────────────────────────────────────────────────
\paragraph{Step 2: Compute residuals using \emph{all} $m$ moments.}

Plug $\hat\beta_r$ into the full set of moment conditions to get the
residual vector
\begin{equation}
  \hat g \;=\; \Delta - Q_H\,\hat\beta_r \;\in\; \mathbb{R}^m.
  \label{eq:residuals}
\end{equation}
The first $k$ elements of $\hat g$ are zero by construction
(from~\eqref{eq:restricted_est}); the remaining $m - k$ elements
are the overidentifying residuals we want to test.

% ── Step 3 ──────────────────────────────────────────────────────────────────
\paragraph{Step 3: Estimate the covariance matrix $\hat\Omega$.}

Estimate $\hat\Omega$ from the data as described in
Section~\ref{sec:gmm}.  This is an $m\times m$ matrix.  It is
rank-deficient (rank $\leq G\times T \ll m$) --- that is the root cause
of Problem 1.  We do \emph{not} try to invert it directly.

% ── Step 4 ──────────────────────────────────────────────────────────────────
\paragraph{Step 4: Form a regularised weighting matrix.}

Rather than inverting $\hat\Omega$ (which we cannot do), add a small
positive constant $\varepsilon$ to every diagonal element and invert
the result:
\begin{equation}
  W \;=\; \bigl(\hat\Omega + \varepsilon I_m\bigr)^{-1},
  \label{eq:reg_inv}
\end{equation}
where $I_m$ is the $m\times m$ identity matrix and $\varepsilon > 0$
is a small regularisation parameter (e.g.\ $\varepsilon = 10^{-6}$).

\medskip
\noindent\textbf{Why this addresses Problem 1.}
Adding $\varepsilon$ to the diagonal nudges all eigenvalues of
$\hat\Omega$ away from zero, making the matrix positive definite and
invertible.  Directions where $\hat\Omega$ has near-zero eigenvalues
(the ``redundant'' moment directions from Section~\ref{subsec:singularity})
end up receiving weight approximately $1/\varepsilon$ in $W$, which is
large.  We will correct for this distortion in Step 6.

\medskip
\noindent Compute the raw test statistic
\begin{equation}
  T \;=\; \hat g'\,W\,\hat g
        \;=\; \hat g'\bigl(\hat\Omega + \varepsilon I_m\bigr)^{-1}\hat g.
  \label{eq:raw_stat}
\end{equation}

% ── Step 5 ──────────────────────────────────────────────────────────────────
\paragraph{Step 5: Understand the distribution of $T$.}

If $\hat\Omega$ were the \emph{true} $\Omega$ and $m$ moments were all
independent, $T$ would follow $\chi^2(m-k)$ under the null.  Neither
condition holds here.

To see what $T$ actually follows, note that under the null hypothesis
$T$ is a \emph{weighted sum of independent chi-squared(1) variables}:
\begin{equation}
  T \;\overset{d}{\approx}\;
    \sum_{i=1}^{m} \eta_i\, Z_i^2,
    \qquad Z_i \overset{iid}{\sim} N(0,1),
  \label{eq:weighted_chisq}
\end{equation}
where $\eta_1,\ldots,\eta_m$ are the eigenvalues of the matrix
\begin{equation}
  A \;:=\; W\hat\Omega
          \;=\; \bigl(\hat\Omega + \varepsilon I_m\bigr)^{-1}\hat\Omega.
  \label{eq:A_matrix}
\end{equation}

\medskip
\begin{tcolorbox}[colback=gray!8, colframe=gray!40,
  title=What are these eigenvalues $\eta_i$?]
If $\hat\Omega$ has eigenvalues $\lambda_1 \geq \lambda_2 \geq \cdots \geq \lambda_m \geq 0$,
then $A = (\hat\Omega + \varepsilon I)^{-1}\hat\Omega$ has eigenvalues
\[
  \eta_i \;=\; \frac{\lambda_i}{\lambda_i + \varepsilon}.
\]
\begin{itemize}[noitemsep]
  \item Directions with \textbf{large} $\lambda_i$ (genuine signal):
        $\eta_i \approx 1$.
  \item Directions with \textbf{near-zero} $\lambda_i$ (redundant
        moments, the null space of $\hat\Omega$): $\eta_i \approx 0$.
\end{itemize}
So $A$ acts like a smooth projection: it assigns nearly full weight
to informative directions and nearly zero weight to redundant ones.
\end{tcolorbox}

\medskip
\noindent Because the $\eta_i$ are not all equal, $T$ is \emph{not}
chi-squared for any standard degrees of freedom.  This is where the
Satterthwaite--Box approximation enters.

% ── Step 6 ──────────────────────────────────────────────────────────────────
\paragraph{Step 6: Apply the Satterthwaite--Box approximation.}

The idea is simple: approximate the distribution of $T$ by a
\emph{scaled} chi-squared, $T \approx c\cdot\chi^2(\nu)$, where the
scale $c$ and degrees of freedom $\nu$ are chosen to match the first
two moments of $T$.

\subparagraph{First moment (mean) of $T$.}
\begin{equation}
  \mathbb{E}[T]
    \;=\; \sum_{i=1}^{m}\eta_i
    \;=\; \operatorname{tr}(A)
    \;=\; \operatorname{tr}\!\bigl((\hat\Omega+\varepsilon I)^{-1}\hat\Omega\bigr).
  \label{eq:ET}
\end{equation}
\textit{In words}: the expected value of $T$ equals the sum of the
effective weights.  If all moments were independent ($\hat\Omega$
diagonal), this would just equal $m - k$ --- the classical degrees
of freedom.  With redundant moments, it is much smaller.

\subparagraph{Second moment (variance) of $T$.}
\begin{equation}
  \operatorname{Var}(T)
    \;=\; 2\sum_{i=1}^{m}\eta_i^2
    \;=\; 2\operatorname{tr}(A^2)
    \;=\; 2\operatorname{tr}\!\bigl((\hat\Omega+\varepsilon I)^{-1}\hat\Omega\bigr)^2.
  \label{eq:VarT}
\end{equation}

\subparagraph{Matching moments.}
Set $c$ and $\nu$ so that $\mathbb{E}[c\chi^2(\nu)] = \mathbb{E}[T]$
and $\operatorname{Var}[c\chi^2(\nu)] = \operatorname{Var}(T)$:
\begin{align}
  c\,\nu       &= \operatorname{tr}(A),                  \label{eq:match1}\\
  2c^2\,\nu    &= 2\operatorname{tr}(A^2).                \label{eq:match2}
\end{align}
Dividing~\eqref{eq:match2} by~\eqref{eq:match1} gives $c$; then $\nu$
follows:
\begin{equation}
  \boxed{
    c \;=\; \frac{\operatorname{tr}(A^2)}{\operatorname{tr}(A)},
    \qquad
    \nu \;=\; \frac{\bigl[\operatorname{tr}(A)\bigr]^2}{\operatorname{tr}(A^2)}.
  }
  \label{eq:cnu}
\end{equation}
Note that $c\nu = \operatorname{tr}(A) = \mathbb{E}[T]$, so $T/c$ has
mean $\nu$ under the null, consistent with $T/c \sim \chi^2(\nu)$.

\medskip
\begin{tcolorbox}[colback=gray!8, colframe=gray!40,
  title=Intuition behind $c$ and $\nu$]
Think of $\nu$ as the \emph{effective number of independent moments}
after accounting for redundancy.  If all $m$ moments were truly
independent ($\eta_i = 1$ for all $i$), then $\operatorname{tr}(A)=m$,
$\operatorname{tr}(A^2)=m$, so $c=1$ and $\nu=m$ --- exactly the
classical result.  As moments become more redundant (many $\eta_i$
near zero), $\operatorname{tr}(A^2)/\operatorname{tr}(A)$ grows because
variance is dominated by fewer large $\eta_i$, and the effective degrees
of freedom $\nu$ shrinks accordingly.
\end{tcolorbox}

% ── Step 7 ──────────────────────────────────────────────────────────────────
\paragraph{Step 7: Account for estimated parameters.}

The derivation above treats $\beta_0$ as known.  Because we estimated
$k$ parameters in Step~1, we lose $k$ degrees of freedom.  Reduce the
effective degrees of freedom by $k$:
\begin{equation}
  \nu^* \;=\; \nu - k
            \;=\; \frac{\bigl[\operatorname{tr}(A)\bigr]^2}{\operatorname{tr}(A^2)}
              - k.
  \label{eq:nu_corrected}
\end{equation}
The corrected test statistic is
\begin{equation}
  J^* \;=\; \frac{T}{c}
           \;=\; \frac{\hat g'W\hat g \cdot \operatorname{tr}(A)}
                      {\operatorname{tr}(A^2)},
  \label{eq:jstar}
\end{equation}
which is compared to the $\chi^2(\nu^*)$ distribution.

% ── Step 8 ──────────────────────────────────────────────────────────────────
\paragraph{Step 8: Compute the $p$-value and make a decision.}

\begin{equation}
  p\text{-value} \;=\; \Pr\!\bigl(\chi^2(\nu^*) \;>\; J^*\bigr).
  \label{eq:pvalue}
\end{equation}
Reject the null hypothesis of correct specification at significance
level $\alpha$ if $p < \alpha$, or equivalently if
$J^* > \chi^2_\alpha(\nu^*)$ (the upper-$\alpha$ quantile of
$\chi^2(\nu^*)$).

% ─────────────────────────────────────────────────────────────────────────────
\subsection{How Each Step Solves a Specific Problem}
\label{subsec:recap}
% ─────────────────────────────────────────────────────────────────────────────

Table~\ref{tab:correction_summary} summarises how the procedure above
addresses the two failures of the standard $J$ test.

\begin{table}[htbp]
\centering
\caption{Problems with the Standard $J$ Test and Their Corrections}
\label{tab:correction_summary}
\begin{tabular}{p{3.2cm} p{4.5cm} p{4.5cm}}
\toprule
 & \textbf{Singularity of $\hat\Omega$}
 & \textbf{Downward Bias} \\
\midrule
\textit{Root cause}
  & All $m$ DiD moments are built from $\leq G\times T$ cohort-time
    averages, so $\hat\Omega$ has rank $\ll m$.
  & $\hat\beta$ from the full set forces residuals orthogonal to $Q_H$
    in the $\hat\Omega^{-1}$ metric, making $\hat g$ artificially small.\\[6pt]
\textit{Consequence}
  & $\hat\Omega^{-1}$ does not exist; $J$ cannot be computed.
    Even with a pseudo-inverse, the $\chi^2(m-k)$ distribution does
    not apply.
  & $\mathbb{E}[J] < m - k$; the test under-rejects --- the model
    looks better than it is. \\[6pt]
\textit{Fix}
  & Steps 4--6: replace $\hat\Omega^{-1}$ with regularised $W$;
    apply Satterthwaite--Box to obtain corrected scale $c$ and
    effective d.f.\ $\nu$.
  & Step 1: estimate $\hat\beta$ from a restricted (just-identified)
    subset of moments; the overidentifying residuals are then
    unconstrained. \\[6pt]
\textit{Intuition}
  & The effective d.f.\ $\nu = [\mathrm{tr}(A)]^2/\mathrm{tr}(A^2)$
    automatically counts only the independent directions in the moment
    space, discarding redundant ones.
  & By using a held-out set of moments for testing, we stop
    ``moving the board after seeing the darts.'' \\
\bottomrule
\end{tabular}
\end{table}

% ─────────────────────────────────────────────────────────────────────────────
\subsection{Complete Algorithm}
\label{subsec:algorithm}
% ─────────────────────────────────────────────────────────────────────────────

For convenience, we collect the full corrected procedure below.

\begin{enumerate}[label=\textbf{Step \arabic*.}, leftmargin=*, itemsep=4pt]

  \item \textbf{Restricted estimation.}
        Choose $k$ linearly independent moment conditions (one per
        parameter).  Solve $\Delta_{\mathrm{res}} =
        Q_{H,\mathrm{res}}\hat\beta_r$ to obtain $\hat\beta_r$.
        \emph{(Fixes downward bias.)}

  \item \textbf{Full residuals.}
        Compute $\hat g = \Delta - Q_H\hat\beta_r \in \mathbb{R}^m$.

  \item \textbf{Covariance estimation.}
        Estimate $\hat\Omega$ using the autocovariance / $R\Omega_v R'$
        method described in Section~\ref{sec:gmm}.

  \item \textbf{Regularised weighting.}
        Set $W = (\hat\Omega + \varepsilon I_m)^{-1}$.
        \emph{(Makes inversion feasible; fixes singularity.)}

  \item \textbf{Raw test statistic.}
        Compute $T = \hat g' W \hat g$.

  \item \textbf{Effective-weight matrix.}
        Compute $A = W\hat\Omega = (\hat\Omega+\varepsilon I)^{-1}\hat\Omega$.

  \item \textbf{Satterthwaite--Box parameters.}
        \[
          c \;=\; \frac{\operatorname{tr}(A^2)}{\operatorname{tr}(A)},
          \qquad
          \nu^* \;=\; \frac{[\operatorname{tr}(A)]^2}{\operatorname{tr}(A^2)} - k.
        \]
        \emph{($\nu^*$ is the effective degrees of freedom after
        accounting for moment redundancy and parameter estimation.)}

  \item \textbf{Corrected test statistic.}
        \[
          J^* \;=\; T/c.
        \]

  \item \textbf{Inference.}
        Reject $H_0$ (correct specification) at level $\alpha$ if
        \[
          J^* \;>\; \chi^2_\alpha(\nu^*),
          \qquad
          \text{or equivalently}\quad
          p = \Pr(\chi^2(\nu^*)>J^*) < \alpha.
        \]

\end{enumerate}

\medskip
\begin{tcolorbox}[colback=blue!5, colframe=blue!30,
  title=Practical note on $\varepsilon$]
The regularisation parameter $\varepsilon$ should be small relative to
the typical non-zero eigenvalues of $\hat\Omega$ but large relative to
numerical noise.  A common choice is $\varepsilon = 10^{-6}$ times
the maximum eigenvalue of $\hat\Omega$.  Sensitivity of $J^*$ and
$\nu^*$ to $\varepsilon$ can be checked by repeating with
$\varepsilon \in \{10^{-4}, 10^{-6}, 10^{-8}\}$; stable results across
this range indicate the correction is reliable.
\end{tcolorbox}
